%%=============================================================================
%% Verbeteringen
%%=============================================================================

\chapter{Verbeteringen}%
\label{ch:verbeteringen}

Deze bijlage bevat de verbeteringen die zijn aangebracht in vergelijking met de scriptie die werd ingediend tijdens de tweede examenperiode.

\section{Herziening en Nieuwe Aanpak}

De scriptie is nagenoeg volledig herzien en herschreven. Hoewel enkele delen van de inleiding zijn hergebruikt, is het merendeel van het document vanaf een nulpunt benaderd, alsof er nog geen kennis was opgedaan. Deze ``zero knowledge'' benadering was cruciaal om de obstakels die in de eerste versie werden aangetroffen, te overtreffen en een nieuwe aanpak te ontwikkelen voor zowel het schrijven van deze scriptie als het opzetten van de proof of concept.

\section{Belangrijkste Veranderingen}

De belangrijkste veranderingen omvatten:
\begin{itemize}
    \item \textbf{Stand van zaken:} Een volledig nieuwe literatuurstudie is uitgevoerd, waarbij relevante bronnen zijn geïdentificeerd en geanalyseerd.
    \item \textbf{Methodologie:} De methodiek is herzien en aangepast om beter aan te sluiten bij de onderzoeksdoelstellingen en om praktische obstakels te omzeilen die eerder werden ervaren.
    \item \textbf{Proof of Concept:} De proof of concept is volledig herontwikkeld met een nieuwe benadering en strategie, waarbij rekening is gehouden met de beperkingen en lessen uit de eerste versie.
    \item \textbf{Conclusie:} De conclusie is opnieuw geformuleerd om de nieuwe bevindingen en inzichten te weerspiegelen die voortkwamen uit de herziene aanpak.
    \item \textbf{Bronnen:} Alle bronnen zijn opnieuw onderzocht en geselecteerd vanaf het punt van de nieuwe start, om ervoor te zorgen dat de gebruikte referenties betrouwbaar zijn, en om te vermijden dat moeilijk vindbare bronnen opnieuw werden gebruikt.
\end{itemize}

Deze herzieningen hebben geleid tot een meer samenhangende scriptie, waarbij de uitdagingen uit de eerste versie zijn overwonnen en een sterkre basis is gelegd voor zowel het academische als het praktische deel van het onderzoek.